\section{Excitation-Coherent Telemetry Analysis}
\label{sec:method}

ECTA processes the tool's telemetry in windows (30~s for full-bus
captures, 60~s for sampling loggers, 120~s or one trip for polling) and
computes three feature families. All features are \emph{self-referenced}:
none requires a healthy recording of the same vehicle or tool.

\subsection{Loss Features}
For a passive capture, each identifier $k$ whose inter-arrival times are
predominantly within $[0.5,1.5]$ of their median $P_k$ is treated as
periodic. A gap $\Delta$ implies $\mathrm{round}(\Delta/P_k)-1$ missing
frames. We aggregate the missing-frame ratio, missing frames and gap events
per minute, all-identifier silences, the longest silence and the frame
rate. For a polling tester the corresponding quantities are derived from the
inter-response gaps relative to the loop period: the share and rate of
timeout-like gaps, the excess time fraction, long gaps ($>0.8$~s, the
re-initialisation signature), the longest gap and the per-PID cycle
completeness.

\subsection{Structure Features}
These describe \emph{how} losses are distributed, which is what separates
the origins of Section~\ref{sec:model}:
\emph{selectivity}---a $\chi^2$ statistic and a Gini coefficient of
per-identifier losses against the airtime-proportional pattern that a DLC
fault produces (an ECU-side fault concentrates losses on few identifiers);
\emph{lateness and bunching}---the share of intervals slightly above the
period without a missing frame, and below half the period (retransmission
and backlog release); \emph{jitter}; and \emph{load coupling}---the local
frame density at the moment of loss relative to the window mean, and the
rank correlation between per-second losses and per-second load (logger
overflow). For polling, the median and dispersion of timeout-gap durations
distinguish a tester re-initialisation from a run of ECU timeouts.

\subsection{Excitation-Coherence Features}
The tool decodes engine and vehicle speed from the frames or responses it
did receive (for passive captures, from broadcast signals; for polling,
from PIDs \texttt{0x0C} and \texttt{0x0D}), forming the observed proxy
$\hat v(t)=w_e n(t)/2000+w_r s(t)/60$ on a 1--2~s grid. Let $y_j$ be the
count of loss events in bin $j$. ECTA fits the Poisson generalized linear
model~\cite{mccullagh1989}
\begin{equation}
y_j\sim\mathrm{Poisson}\big(\exp(\alpha+\beta\,x_j)\big),\quad
x_j=\mathrm{std}\big(\ln(\hat v_j+0.05)\big),
\label{eq:glm}
\end{equation}
by iteratively reweighted least squares and records the slope
$\hat\beta$, its Wald statistic $z_\beta=\hat\beta/\widehat{\mathrm{se}}(\hat\beta)$,
the Spearman correlation between $y$ and $\hat v$, and the log ratio of
loss rates above and below the median excitation. Under (\ref{eq:rate}),
DLC losses have $\beta>0$; overflow and random loss have $\beta\approx0$.
The statistics are computed twice: for all loss events, and for strong
events only (silences or gaps exceeding three loop periods), which isolates
brown-outs from the tool's own background irregularity.

On its own, $z_\beta$ is a training-free test: it needs no labels and no
healthy data from the target vehicle. We report it as a separate method.

\subsection{Detection and Attribution}
The feature vector (21 passive, 21 polling features) is classified by
gradient-boosted decision trees (LightGBM~\cite{lightgbm}; 400 trees,
learning rate 0.05, 31 leaves, class-balanced weights). The same features
feed a binary detector (wear versus everything else, including
confounders) and a five-way attribution classifier (healthy, DLC wear,
ECU-side, bus EMI, overflow/random loss). The model is deliberately
small (a 1.3-MB tree ensemble) and the feature extraction is a single pass
over the window, so ECTA can run on the tool or at the fleet back end.

\subsection{Health Index and Remaining Useful Life}
Because the fault is persistent, evidence accumulates across trips. For
prognosis, a gradient-boosted regressor maps per-trip features to an
estimate $h_k$ of $\log_{10}R$ (the health index, HI). Two uses of the HI
are evaluated.

\emph{Trend alarm.} A least-squares line through the last 30 trips' HI
values is extrapolated in calendar time to the failure level; an alarm is
raised when the projected crossing is less than 30 days away.

\emph{Bayesian exposure-domain RUL.} Here the wear law (\ref{eq:wear}) is
fitted in the \emph{exposure} domain, where degradation is physically
stationary; each trip's exposure $N_k$ is computed from its own telemetry
via (\ref{eq:exposure}). With $R_0$, $R_{\max}$ and $w/N_{50}$ fixed, the
unit-specific $\ln N_{50}$ receives a Gaussian prior estimated from the
failure exposures of other vehicles. HI readings above a detection floor
$h_0$ are Gaussian observations of $\log_{10}R(N_i)$ with noise $\tau$;
readings at or below $h_0$ are treated as left-censored, contributing
$\Phi\big((h_0-\log_{10}R(N_i))/\tau\big)$ to the likelihood, because a
low-rate tool cannot distinguish a fresh contact from a moderately worn
one. Survival to the current exposure truncates the posterior. The RUL in
days is the posterior median of $(N_f-N_k)/\dot N_k$, where $N_f$ solves
$R(N_f)=R_{\mathrm{fail}}$ and $\dot N_k$ is the vehicle's mean exposure
per day so far, so heavily used vehicles are predicted to fail sooner.
The floor $h_0$ (99th percentile of the HI on contacts below
1\,$\Omega$), $\tau$ and the prior are calibrated on one half of the
vehicles and applied to the other.
